\section{Experiments}
\label{sec:experiments}
\subsection{Models and Training Details}
We employ the Qwen3-4B-Base~\citep{Yang2025Qwen3TR} and Qwen3-8B-Base models to assess the impact of scale within a single architectural family. Second, to ensure our approach is effective on a distinct lineage, we utilize the OctoThinker-3B and OctoThinker-8B models~\citep{wang2025octothinker}.This choice is particularly relevant as \cite{wang2025octothinker} reported that applying RL training directly to Llama models yielded suboptimal results. As the OctoThinker series is continually trained from the Llama-3.1 models~\citep{Dubey2024TheL3}, this comprehensive selection allows us to test our framework across different foundational architectures -- Qwen vs. Llama. We assess our framework on a comprehensive suite of benchmarks, with the evaluation benchmarks presented in Appendix~\ref{app:tasks_and_evaluation}.

Our entire framework is implemented based on the EasyR1 codebase~\citep{zheng2025easyr1}. In each iteration of the \method{} co-evolutionary loop, we follow a specific set of hyperparameters. The Challenger ($Q_{\theta}$) first generates a candidate pool of $N=8,000$ questions. To construct the training dataset for the Solver, these questions are filtered based on consistency. For each candidate question, we sample $m=10$ answers from the current Solver ($S_{\phi}$). 
A question is retained for the training set only if the number of answers matching the majority-vote pseudo-label is between 3 and 7, inclusive ($\delta=0.25$). This numerical range is consistent with the methodology used in previous research~\citep{Zhang2025GRPOLEADAD,li2025semanticallyawarerewardsopenendedr1,Bercovich2025LlamaNemotronER}.
When training the Challenger, the uncertainty reward $r(x; \phi)$ is calculated by sampling $m=10$ responses from the Solver. For the intra-batch repetition penalty, we set the clustering distance threshold to $\tau_{\text{BLEU}} = 0.5$. 


\subsubsection{Evaluation Setting}

The evaluation code is adopted from General-Reasoner~\citep{Ma2025GeneralReasonerAL}. To ensure consistency, we reran the released evaluation code and reported the corresponding results. The reproduced results are well aligned with General-Reasoner and those in the Qwen-3 technical report~\citep{Yang2025Qwen3TR}.

The results presented in all experimental tables are obtained after 45 training steps, while in the figures we report evaluations performed at checkpoints every 15 steps during solver training. All results are reported based on the held-out test sets, ensuring a fair comparison across baselines and reproduced methods.
Further implementation details and  prompts  can be found in Appendix~\ref{app:traininghyper}.

\subsection{Results in Mathematical Reasoning}
The comprehensive results of our experiments are presented in Table~\ref{tab:final_main_results}. The findings confirm that our proposed framework, \method{}, is a highly effective, model-agnostic method for enhancing the performance of language models on mathematical tasks across different architectures and scales.

\begin{table*}[t]
\centering
\small
\caption{Comprehensive results on mathematical reasoning benchmarks. We compare each base model against a \method{} (\raisebox{-0.2ex}{\includegraphics[width=0.02\textwidth]{figures/snowflake.png}} challenger) baseline (where the Solver is trained on questions from an untrained Challenger) and our method, \method{}. The peak performance  is highlighted in \textbf{bold}.}
% \vspace{-0.1in}
\label{tab:final_main_results}
\scalebox{0.9}{
\begin{tabular}{@{}l>{\columncolor{gray!15}}cccccccc@{}}
\toprule
\textbf{Model Name} & \textbf{Avg.} & \textbf{AMC} & \textbf{Minerva} & \textbf{MATH} & \textbf{GSM8K} & \textbf{Olympiad} & \textbf{AIME25} & \textbf{AIME24}\\
\midrule
\multicolumn{8}{@{}l}{\textit{Qwen3-4B-Base}} \\
\quad Base Model (w/o training)     & 42.57 & 45.70 & 38.24 & 68.20 & 87.79 & 41.04 & 10.30 & 6.7\\
\quad Absolute Zero & 46.42 & 52.45 & 41.96 & 76.20 & 89.34 & 42.56 & 10.20 & 12.20 \\
\quad \method{} (\raisebox{-0.3ex}{\includegraphics[width=0.02\textwidth]{figures/snowflake.png}} challenger)  & 45.01 & 45.00 & 45.22 & 72.80 & 87.87 & 41.19 & \textbf{10.20} & 12.80\\
\textbf{\quad \method{} (our method)} & \textbf{49.93} & \textbf{57.27} & \textbf{52.94} & \textbf{79.60} & \textbf{92.12} & \textbf{44.59} & 9.60 & \textbf{13.40}\\
\midrule
\multicolumn{8}{@{}l}{\textit{Qwen3-8B-Base}} \\
\quad Base Model (w/o training)  & 48.64 & 51.95 & 50.00 & 78.00 & 89.08 & 44.74 & 12.10 & 14.60 \\
\quad Absolute Zero & 52.68 & 57.89 & 57.90 & 76.60 & 92.00 & 47.80 & 18.20 & 18.40 \\
\quad \method{} (\raisebox{-0.3ex}{\includegraphics[width=0.02\textwidth]{figures/snowflake.png}} challenger)  & 52.10     & 60.70     & 57.72     & 81.60     & 92.56     & 46.44     & 11.60     & 14.10      \\
\textbf{\quad \method{} (our method)} & \textbf{53.72} & \textbf{61.67} & \textbf{60.66} & \textbf{82.00} & \textbf{94.09} & \textbf{48.89} & \textbf{13.30} & \textbf{15.40} \\
\midrule\midrule
\multicolumn{8}{@{}l}{\textit{OctoThinker-3B}} \\
\quad Base Model (w/o training) & 26.64 & 17.19 & 24.26 & \textbf{55.00} & 73.69 & 16.15 & 0.21 & 0.00 \\
\quad Absolute Zero & 27.23 & 22.50 & 22.70 & 53.20 & 75.80 & 13.20 & 0.50 & 2.70 \\
\quad \method{} (\raisebox{-0.3ex}{\includegraphics[width=0.02\textwidth]{figures/snowflake.png}} challenger)  & 27.51     & 20.19     & 24.63     & 54.60     & \textbf{74.98}     & 15.70     & 0.10     & \textbf{2.40}     \\
\textbf{\quad \method{} (our method)} & \textbf{29.32} & \textbf{27.03} & \textbf{27.57} & 54.20 & \textbf{74.98} & \textbf{18.22} & \textbf{3.23} & 0.00 \\
\midrule
\multicolumn{8}{@{}l}{\textit{OctoThinker-8B}} \\
\quad Base Model (w/o training)  & 36.41 & 32.11 & 41.91 & 65.20 & 86.96 & 26.52 & \textbf{1.56} & 0.62 \\
\quad Absolute Zero & 36.60 & 32.50 & 44.90 & 62.80 & 87.00 & 25.60 & 3.30 & 0.10 \\
\quad \method{} (\raisebox{-0.3ex}{\includegraphics[width=0.02\textwidth]{figures/snowflake.png}} challenger)  & 36.98 & 29.30 & 42.28 & 66.20 & \textbf{88.10} & \textbf{27.56} & 1.04 & \textbf{4.38} \\
\textbf{\quad \method{} (our method)} & \textbf{38.52} & \textbf{34.03} & \textbf{48.22} & \textbf{68.80} & 87.19 & \textbf{27.56} & 0.42 & 3.44 \\
\bottomrule
\end{tabular}}
\end{table*}


Our iterative training process consistently and substantially improves upon the performance of the base models. This holds true for large models like Qwen3-8B-Base, where three iterations of \method{} raise the average performance from a baseline of 49.18 to 54.69, a significant gain of \textbf{+5.51} points. Similarly, on the smaller OctoThinker-3B, our method improves the average score from 26.64 to 29.32 (\textbf{+2.68} points), demonstrating the broad applicability of our self-supervised training loop.

% This improvement is progressive, with the results showing a clear trend of performance gains across iterations. For instance, the Qwen3-8B-Base model's average score climbs from a base performance of 49.18 to 53.39 (Iter 1) and ultimately reaches 54.69 (Iter 3). A similar monotonic improvement is observed on OctoThinker-3B, which progresses from its base score of 26.64 to 29.32 after three iterations. This consistent growth underscores the benefits of the co-evolutionary dynamic, where the progressively more capable Solver learns from an increasingly challenging curriculum.

The critical role of the Challenger's RL-based training is validated by the immediate performance leap from the Base Challenger to the first iteration of \method{}. On Qwen3-8B-Base, this first iteration provides a +1.52 point gain over the baseline, and the improvement is even more pronounced on Qwen3-4B-Base at +3.7 points. This confirms that the intelligent curriculum generated by the RL-trained Challenger is significantly more effective than that of a non-trained generator.

\subsection{Results in General Reasoning}

\begin{table*}[t]
\centering
\small
\setlength{\tabcolsep}{3mm}{
\caption{Results on general-domain reasoning benchmarks. The table compares the Base Model, a \method{} (\raisebox{-0.2ex}{\includegraphics[width=0.02\textwidth]{figures/snowflake.png}} challenger) baseline, and our  \method{}. The peak performance is highlighted in \textbf{bold}.}
\label{tab:general_domain_results}
% \vspace{-0.1in}
\scalebox{0.92}{
\begin{tabular}{@{}l>{\columncolor{gray!15}}cccc@{}}
\toprule
\textbf{Model Name} & \textbf{Overall Avg.} &  \textbf{SuperGPQA} & \textbf{MMLU-Pro} & \textbf{BBEH} \\
\midrule
\multicolumn{5}{@{}l}{\textit{Qwen3-4B-Base}} \\
\quad Base Model (w/o training) & 26.34 & 20.88 & 50.58 & 7.57 \\
\quad Absolute Zero & 29.33 & 27.10 & 52.60 & 8.30 \\

\quad \method{} (\raisebox{-0.3ex}{\includegraphics[width=0.02\textwidth]{figures/snowflake.png}} challenger) & 28.52 & 24.77 & 54.20 & 6.59 \\
\textbf{\quad \method{} (our method)} & \textbf{31.15} & \textbf{27.55} & \textbf{55.47} & \textbf{10.42} \\
\midrule[0.5pt]
\multicolumn{5}{@{}l}{\textit{Qwen3-8B-Base}} \\
\quad Base Model (w/o training) & 31.98 & 28.33 & 58.97 & 8.63 \\
\quad Absolute Zero & 34.40 & 31.89 & 60.50 & 10.80 \\

\quad \method{} (\raisebox{-0.3ex}{\includegraphics[width=0.02\textwidth]{figures/snowflake.png}} challenger) & 31.29       &  30.12       & 54.14       & 9.60       \\

\textbf{\quad \method{} (our method)} & \textbf{34.50} & \textbf{31.38} & \textbf{61.53} & \textbf{10.60} \\
\midrule\midrule
\multicolumn{5}{@{}l}{\textit{OctoThinker-3B}} \\
\quad Base Model (w/o training) & 7.47 & 10.09 & 10.87 & 1.46 \\
\quad Absolute Zero & 16.03 & 17.70 & 24.30 & 6.10 \\
\quad \method{} (\raisebox{-0.3ex}{\includegraphics[width=0.02\textwidth]{figures/snowflake.png}} challenger) & 10.04   & 11.19       & 14.53       & \textbf{4.40}       \\
\textbf{\quad \method{} (our method)} & \textbf{11.12}  & \textbf{12.44} & \textbf{16.71} & 4.20 \\
\midrule[0.5pt]
\multicolumn{5}{@{}l}{\textit{OctoThinker-8B}} \\
\quad Base Model (w/o training) & 11.70 & 13.26 & 20.21 & 1.64 \\
\quad Absolute Zero & 18.40 & 18.80 & 31.40 & 5.00 \\
\quad \method{} (\raisebox{-0.3ex}{\includegraphics[width=0.02\textwidth]{figures/snowflake.png}} challenger) & 21.30  & 16.99 & \textbf{41.46} & 5.46 \\
\textbf{\quad \method{} (our method)} & \textbf{23.00} & \textbf{19.82} & 40.92 & \textbf{8.25} \\
\bottomrule
\end{tabular}}}
\end{table*}
Previous work has demonstrated that training language models on  reasoning-intensive domains, such as mathematics, can lead to improvements in general-domain capabilities~\citep{Huan2025DoesMR}. A key question, however, is whether this generalization effect still holds when the training curriculum is not human-labeled, but entirely self-generated through \method{}. 



As shown in Table~\ref{tab:general_domain_results}, this transfer of skills is evident across all tested models. For instance, three iterations of our math-focused training improve the average general-domain score of Qwen3-8B-Base by +5.13 points and OctoThinker-3B by +3.65 points. This generalization also extends to the key performance patterns observed in the mathematical results. This confirms that our method does not merely teach domain-specific knowledge, but enhances the model's underlying capabilities in a way that successfully generalizes across domains.
